%% =================================================================
%% Kingma, Ba. "Adam: A Method for Stochastic Optimization."
%% ICLR 2015. arXiv:1412.6980v9.
%%
%% Reconstruction of page 4 (printed p. 4) as study notes.
%% Generated by: reproduce-all-pages-basic-tex (batch 1-15)
%%
%% Structure: Section 4 CONVERGENCE ANALYSIS (regret definition eq 5,
%% Theorem 4.1 with unnumbered bound, Corollary 4.2) + start of
%% Section 5 RELATED WORK (cut off at page end).
%% New numbered equation: (5). Equation counter after page: 5.
%% =================================================================

\documentclass[11pt]{article}

\usepackage[margin=1in]{geometry}
\usepackage{amsmath, amssymb}
\usepackage{mathrsfs}
\usepackage{bm}
\usepackage{xcolor}
\usepackage{fancyhdr}

\pagestyle{fancy}
\fancyhf{}
\lhead{\small Published as a conference paper at ICLR 2015}
\rhead{}
\cfoot{\small 4 $\to$ \arabic{page}}
\renewcommand{\headrulewidth}{0.4pt}

\setcounter{page}{1}
\setcounter{equation}{4}

\begin{document}

\section*{4 \textsc{Convergence analysis}}

We analyze the convergence of Adam using the online learning
framework proposed in (Zinkevich, 2003). Given an arbitrary,
unknown sequence of convex cost functions
$f_1(\theta), f_2(\theta), \ldots, f_T(\theta)$. At each time $t$,
our goal is to predict the parameter $\theta_t$ and evaluate it on
a previously unknown cost function $f_t$. Since the nature of the
sequence is unknown in advance, we evaluate our algorithm using
the \emph{regret}, that is the sum of all the previous difference
between the online prediction $f_t(\theta_t)$ and the best fixed
point parameter $f_t(\theta^*)$ from a feasible set $\mathcal{X}$
for all the previous steps. Concretely, the regret is defined as:
%
\begin{equation}
R(T) = \sum_{t=1}^{T} [f_t(\theta_t) - f_t(\theta^*)]
\end{equation}
%
where $\theta^* = \arg\min_{\theta \in \mathcal{X}} \sum_{t=1}^{T} f_t(\theta)$.
We show Adam has $O(\sqrt{T})$ regret bound and a proof is given
in the appendix. Our result is comparable to the best known bound
for this general convex online learning problem. We also use some
definitions simplify our notation, where
$g_t \triangleq \nabla f_t(\theta_t)$ and $g_{t,i}$ as the
$i^\text{th}$ element. We define
$g_{1:t,i} \in \mathbb{R}^t$ as a vector that contains the
$i^\text{th}$ dimension of the gradients over all iterations till
$t$, $g_{1:t,i} = [g_{1,i}, g_{2,i}, \ldots, g_{t,i}]$. Also, we
define $\gamma \triangleq \frac{\beta_1^2}{\sqrt{\beta_2}}$. Our
following theorem holds when the learning rate $\alpha_t$ is
decaying at a rate of $t^{-\frac{1}{2}}$ and first moment running
average coefficient $\beta_{1,t}$ decay exponentially with
$\lambda$, that is typically close to 1, e.g.\
$1 - 10^{-8}$.

\vspace{0.5em}

\noindent\textbf{Theorem 4.1.} \emph{Assume that the function
$f_t$ has bounded gradients,
$\|\nabla f_t(\theta)\|_2 \leq G$,
$\|\nabla f_t(\theta)\|_\infty \leq G_\infty$ for all
$\theta \in \mathbb{R}^d$ and distance between any $\theta_t$
generated by Adam is bounded,
$\|\theta_n - \theta_m\|_2 \leq D$,
$\|\theta_m - \theta_n\|_\infty \leq D_\infty$ for any
$m, n \in \{1, \ldots, T\}$, and $\beta_1, \beta_2 \in [0, 1)$
satisfy $\frac{\beta_1^2}{\sqrt{\beta_2}} < 1$. Let
$\alpha_t = \frac{\alpha}{\sqrt{t}}$ and
$\beta_{1,t} = \beta_1 \lambda^{t-1}$, $\lambda \in (0, 1)$.
Adam achieves the following guarantee, for all $T \geq 1$.}
%
\[
R(T) \leq \frac{D^2}{2\alpha(1 - \beta_1)}
          \sum_{i=1}^{d} \sqrt{T \, \widehat{v}_{T,i}}
        + \frac{\alpha(1 + \beta_1) G_\infty}
               {(1 - \beta_1)\sqrt{1 - \beta_2}\,(1 - \gamma)^2}
          \sum_{i=1}^{d} \|g_{1:T,i}\|_2
        + \sum_{i=1}^{d}
          \frac{D_\infty^2 \, G_\infty \sqrt{1 - \beta_2}}
               {2\alpha(1 - \beta_1)(1 - \lambda)^2}
\]

\noindent
Our Theorem~4.1 implies when the data features are sparse and
bounded gradients, the summation term can be much smaller than
its upper bound
$\sum_{i=1}^{d} \|g_{1:T,i}\|_2 \ll dG_\infty\sqrt{T}$ and
$\sum_{i=1}^{d} \sqrt{T\,\widehat{v}_{T,i}} \ll dG_\infty\sqrt{T}$,
in particular if the class of function and data features are in
the form of section~1.2 in (Duchi et al., 2011). Their results
for the expected value
$\mathbb{E}[\sum_{i=1}^{d} \|g_{1:T,i}\|_2]$ also apply to Adam.
In particular, the adaptive method, such as Adam and Adagrad, can
achieve $O(\log d\sqrt{T})$, an improvement over $O(\sqrt{dT})$
for the non-adaptive method. Decoupling $\beta_1$ towards zero is
important in our theoretical analysis and also matches previous
empirical findings, e.g.\ (Sutskever et al., 2013) suggests
reducing the momentum coefficient in the end of training can
improve convergence.

Finally, we can show the average regret of Adam converges,

\vspace{0.5em}

\noindent\textbf{Corollary 4.2.} \emph{Assume that the function
$f_t$ has bounded gradients,
$\|\nabla f_t(\theta)\|_2 \leq G$,
$\|\nabla f_t(\theta)\|_\infty \leq G_\infty$ for all
$\theta \in \mathbb{R}^d$ and distance between any $\theta_t$
generated by Adam is bounded,
$\|\theta_n - \theta_m\|_2 \leq D$,
$\|\theta_m - \theta_n\|_\infty \leq D_\infty$ for any
$m, n \in \{1, \ldots, T\}$. Adam achieves the following
guarantee, for all $T \geq 1$.}
%
\[
\frac{R(T)}{T} = O\!\left(\frac{1}{\sqrt{T}}\right)
\]

\noindent
This result can be obtained by using Theorem~4.1 and
$\sum_{i=1}^{d} \|g_{1:T,i}\|_2 \leq dG_\infty\sqrt{T}$. Thus,
$\lim_{T \to \infty} \frac{R(T)}{T} = 0$.

\vspace{0.5em}

\section*{5 \textsc{Related work}}

Optimization methods having a direct relation to Adam are RMSProp
(Tieleman \& Hinton, 2012; Graves, 2013) and AdaGrad (Duchi et al.,
2011); these relationships are discussed below. Other stochastic
optimization methods include vSGD (Schaul et al., 2012), AdaDelta
(Zeiler, 2012) and the natural Newton method from Roux \&
Fitzgibbon (2010), all setting stepsizes by estimating curvature

% [page ends here — continues on page 5]

\end{document}
